% Review corrections: Analysis 021 observations/O015-review-corrections.md; 10_review_figures.py.
\section{Discussion and conclusion}
\label{sec:conclusion}

Selecting an intervention requires local zero fractions, operation
accounting, validation quality and absolute execution time. Their rankings
can disagree: pressure's value depends on threshold and target set, and
broader attention sparsity need not yield lower latency. Native-relative
speedup also confounds recipe comparison with each reference's cost.

These conclusions are descriptive. The 54 configurations use one training
seed per condition and the same validation split, not 54 independent
replications or an untouched confirmation evaluation. Seven-site OL1 changes
pressure targets and objective normalization together; its norm-cap activity
does not identify why individual activation distributions change. The
$\kappa=0.5$ result is the largest sampled threshold, with no observations
between 0.1 and 0.5. Cross-size comparisons hold tokens, not tokens per
parameter, fixed, use a lower 410M learning rate, and lack larger-model
pressure-free multisite controls. They establish recurrence of one
within-size ordering, not improving quality cost with scale.

The kernel results concern one RTX5090 and batch-size-one, uncached
full-sequence inference. Projection skipping is profitable at selected
settings; attention skipping is not, even where instructions are bypassed.
Stronger compiled dense baselines, cached decoding, other shapes and devices,
independent seeds, and independent confirmation data remain untested.
The study consequently supports conditional intervention choices and
identifies where further controls or implementation work are needed.
